\chapter{Kommentierte Quellen}

Achtundzwanzig Werke, eingeordnet: Zu jedem steht, wofür man es aufschlägt und was es leistet, das andere nicht leisten.

\lit{1}{Andrieu, C., Doucet, A., Holenstein, R.: \textit{Particle Markov Chain Monte Carlo Methods}. J.\ Roy.\ Statist.\ Soc.\ B 72, 2010.}{Der Aufsatz, der Kapitel~7 und~12 verbindet: ein verrauschter Likelihood-Schätzer im Metropolis-Quotienten, und die Kette ist trotzdem exakt. \textbf{Die Grundlage der pseudo-marginalen Verfahren}; schwer, aber die Einleitung ist für jeden lesbar.}

\lit{2}{Asmussen, S., Glynn, P.: \textit{Stochastic Simulation: Algorithms and Analysis}. Springer 2007.}{\textbf{Das mathematisch strengste Lehrbuch des Fachs} --- Varianzreduktion, seltene Ereignisse, Markow-Ketten, stochas\-tische Differen\-tial\-glei\-chungen, mit Beweisen. Die Fundstelle für alles aus den Teilen~II bis IV, was hier nur mit Simulation belegt wurde. Für Leser, die Band 25 beherrschen.}

\lit{3}{Brooks, S., Gelman, A., Jones, G., Meng, X.-L. (Hg.): \textit{Handbook of Markov Chain Monte Carlo}. Chapman \& Hall 2011.}{Das Nachschlagewerk zu Teil~III: Geyers Einführung, Neals Kapitel über Hamilton-Monte-Carlo, Diagnostik, Anwendungen. \textbf{Neals Beitrag ist die beste Darstellung von Kapitel~9} und frei verfügbar.}

\lit{4}{Del Moral, P.: \textit{Feynman--Kac Formulae: Genealogical and Interacting Particle Systems}. Springer 2004.}{Die Theorie, die Teilchenfilter und Aufspaltung als dasselbe Objekt behandelt. \textbf{Nicht zum Lesen, zum Nachschlagen} --- die Fundstelle für die Erwartungstreue in Satz 12.2.}

\lit{5}{Doucet, A., de Freitas, N., Gordon, N. (Hg.): \textit{Sequential Monte Carlo Methods in Practice}. Springer 2001.}{Die Sammlung, die den Teilchenfilter bekannt machte --- Kapitel~12 mit allen Varianten und Anwendungen von der Navigation bis zur Ökonometrie. Der Beitrag von Doucet, Godsill und Andrieu ist die Einführung.}

\lit{6}{Gelman, A., Carlin, J., Stern, H., Dunson, D., Vehtari, A., Rubin, D.: \textit{Bayesian Data Analysis}. 3.\ Aufl., Chapman \& Hall 2013.}{Das Standardwerk der angewandten Bayes-Statistik, mit den Kapiteln~10 bis 12 zu Metropolis, Gibbs, Hamilton und Diagnostik. \textbf{Die Quelle für $\hat R$, die Halbketten-Regel und die Warnung vor dem Ausdünnen} aus Kapitel~8.}

\lit{7}{Giles, M.: \textit{Multilevel Monte Carlo Path Simulation}. Oper.\ Res.\ 56, 2008.}{Satz 9.5 im Original, mit dem Kostenbeweis und den Beispielen aus der Finanzmathematik. \textbf{Zwölf Seiten, die ein Fach verändert haben}; Giles' Übersicht in Acta Numerica 2015 führt weiter.}

\lit{8}{Glasserman, P.: \textit{Monte Carlo Methods in Financial Engineering}. Springer 2004.}{Die Simulation aus der Sicht der Finanzmathematik: Varianzreduktion, Quasi-Monte-Carlo mit Brownscher Brücke, Pfaddiskretisierung, seltene Ereignisse. \textbf{Die beste Darstellung der Kapitel~5, 6 und~9 für Leser mit Band 25}; die Beispiele sind Optionen, die Methoden allgemein.}

\lit{9}{Hammersley, J., Handscomb, D.: \textit{Monte Carlo Methods}. Methuen 1964.}{Die erste Monographie --- antithetische Variablen, Kontrollvariablen, Wichtigkeit, auf $180$ Seiten. \textbf{Noch heute lesbar und ein Dokument des Anfangs}: Vieles, was später wiederentdeckt wurde, steht hier.}

\lit{10}{Hedtstück, U.: \textit{Simulation diskreter Prozesse}. Springer Vieweg 2013.}{\textbf{Auf Deutsch}: die ereignisgesteuerte Simulation von Grund auf --- Ereignisliste, Zustand, Zufallszahlen, Warteschlangen, statistische Auswertung ---, mit Beispielen aus Fertigung und Verkehr. Die deutschsprachige Entsprechung zu Kapitel~10 und~11; knapp und für Leser ohne Band 25 zugänglich.}

\lit{11}{Kloeden, P., Platen, E.: \textit{Numerical Solution of Stochastic Differential Equations}. Springer 1992.}{Das Standardwerk zu Kapitel~9.3: Euler--Maruyama, Milstein, starke und schwache Ordnung mit allen Beweisen, höhere Verfahren. Die Fundstelle für die vierte Auslassung.}

\lit{12}{Knuth, D.: \textit{The Art of Computer Programming, Bd.~2: Seminumerical Algorithms}. 3.\ Aufl., Addison-Wesley 1997.}{Das Kapitel über Zufallszahlen ist die Quelle von Kapitel~2 --- lineare Generatoren, der Satz von Hull und Dobell (erste Auslassung), der Spektraltest, statistische Prüfungen. \textbf{Mit Knuths Geschichte des Generators, den er selbst falsch baute.}}

\lit{13}{Kolonko, M.: \textit{Stochastische Simulation}. Vieweg+Teubner 2008.}{\textbf{Auf Deutsch, mit Beweisen}: Zufallszahlen und ihre Prüfung, Erzeugung von Verteilungen, Varianzreduktion, Markow-Ketten-Verfahren, Auswertung abhängiger Ausgaben. Die Kapitel~2, 4, 5 und~7 dieses Bandes in strengerer Form; die Fundstelle für Leser, die die Beweise in deutscher Sprache lesen wollen.}

\lit{14}{Kroese, D., Taimre, T., Botev, Z.: \textit{Handbook of Monte Carlo Methods}. Wiley 2011.}{Ein Nachschlagewerk mit Programmcode für fast alles in diesem Band: Erzeugung von Verteilungen, Varianzreduktion, Markow-Ketten, seltene Ereignisse mit Kreuzentropie und Aufspaltung. \textbf{Für Leser, die es umsetzen wollen.}}

\lit{15}{L'Ecuyer, P.: \textit{Random Number Generation}. In: Handbook of Computational Statistics, Springer 2012.}{Die moderne Übersicht über Generatoren --- kombinierte, Mersenne-Twister, Zählergeneratoren, Ströme und Teilströme für parallele Läufe. \textbf{Die Quelle für die Regeln in Kapitel~13} zu abgeleiteten Strömen.}

\lit{16}{Law, A.: \textit{Simulation Modeling and Analysis}. 5.\ Aufl., McGraw-Hill 2015.}{Das Standardwerk der ereignisgesteuerten Simulation --- Teil~IV --- und der Statistik der Ausgaben, Kapitel~4. \textbf{Mit dem Kapitel über gemeinsame Zufallszahlen und ihre Synchronisation}, das Aufgabe 4.3 zugrunde liegt. Ingenieurnah, mit vielen Fallstudien.}

\lit{17}{Liu, J.: \textit{Monte Carlo Strategies in Scientific Computing}. Springer 2001.}{Wichtigkeitsstichprobe, Markow-Ketten und Teilchenfilter aus einer Hand, mit den Verbindungen, die dieser Band betont --- der Unabhängigkeitsvorschlag als Gegenstück zur Wichtigkeit, der Filter als Wichtigkeit über die Zeit. \textbf{Das Buch für Leser, die die Einheit des Fachs sehen wollen.}}

\lit{18}{Metropolis, N., Rosenbluth, A., Rosenbluth, M., Teller, A., Teller, E.: \textit{Equation of State Calculations by Fast Computing Machines}. J.\ Chem.\ Phys.\ 21, 1953.}{Der Algorithmus im Original, auf sechs Seiten, mit den harten Scheiben in der Schachtel. \textbf{Lesenswert für die Klarheit} --- und für den Satz, dass das Verfahren \enquote{obviously} die richtige Verteilung liefert; der Beweis ist ein Absatz.}

\lit{19}{Müller-Gronbach, T., Novak, E., Ritter, K.: \textit{Monte Carlo-Algorithmen}. Springer 2012.}{\textbf{Auf Deutsch, mathematisch streng}: Zufallszahlen, Integration, Fluch der Dimension, Varianzreduktion, Markow-Ketten, stochastische Differentialgleichungen; mit Beweisen und mit der Frage, was Zufall gegenüber deterministischen Verfahren prinzipiell leisten kann (Kapitel~3 und~6). Die deutschsprachige Fundstelle für die Sätze dieses Bandes und die natürliche Fortsetzung für Leser, die Band 20 kennen.}

\lit{20}{Neal, R.: \textit{MCMC Using Hamiltonian Dynamics}. In: [3], 2011.}{Die Einführung in Kapitel~9.2 --- Sprungverfahren, Volumenerhaltung, Abstimmung von $\varepsilon$ und $L$, die Klippe des Bruchfalls. Frei verfügbar. \textbf{Wer eine Sache aus diesem Band weiterliest, sollte diese wählen.}}

\lit{21}{Niederreiter, H.: \textit{Random Number Generation and Quasi-Monte Carlo Methods}. SIAM 1992.}{Die Theorie der Diskrepanz, der Satz von Koksma--Hlawka (zweite Auslassung), die Konstruktion der Folgen von Sobol und Halton. \textbf{Zahlentheoretisch streng}; die Quelle für Kapitel~6.4.}

\lit{22}{Owen, A.: \textit{Monte Carlo Theory, Methods and Examples}. Stanford, online, 2013 ff.}{\textbf{Das vollständigste moderne Lehrbuch, frei verfügbar, laufend erweitert}: Zufallszahlen, Varianzreduktion, Quasi-Monte-Carlo mit Verwürfelung, Markow-Ketten, seltene Ereignisse. Owens eigene Erfindungen --- die verwürfelten Netze --- stehen in Kapitel~17. Die Hauptquelle dieses Bandes neben [2] und [24].}

\lit{23}{Robert, C., Casella, G.: \textit{Monte Carlo Statistical Methods}. 2.\ Aufl., Springer 2004.}{Die Markow-Ketten-Verfahren aus Sicht der Statistik, mit Beweisen: Metropolis--Hastings, Gibbs, Konvergenz, Diagnostik, Rao-Blackwellisierung. \textbf{Die Fundstelle für Satz 7.2 in voller Allgemeinheit} und für die Theorie hinter Kapitel~8.}

\lit{24}{Ross, S.: \textit{Simulation}. 6.\ Aufl., Academic Press 2022.}{\textbf{Der Einstieg}: Zufallszahlen, Erzeugung von Verteilungen, ereignisgesteuerte Simulation, Varianzreduktion, Bootstrap --- knapp, mit vielen Aufgaben, ohne mehr Voraussetzungen als Band 15. Die Kapitel~1, 2 und 10 dieses Bandes folgen seiner Anlage.}

\lit{25}{Rubinstein, R., Kroese, D.: \textit{Simulation and the Monte Carlo Method}. 3.\ Aufl., Wiley 2017.}{Der Klassiker von 1981 in dritter Auflage, mit der Kreuzentropie-Methode (Kapitel~11, Forschungskasten) aus der Feder ihres Erfinders. Breit, mit Aufgaben; die Kapitel zu seltenen Ereignissen sind der Grund, es zu besitzen.}

\lit{26}{Saltelli, A., Ratto, M., Andres, T., Campolongo, F., Cariboni, J., Gatelli, D., Saisana, M., Tarantola, S.: \textit{Global Sensitivity Analysis: The Primer}. Wiley 2008.}{Sobol-Indizes, ihre Schätzung aus $N(d+2)$ Läufen, Versuchspläne --- Kapitel~14 in Breite, mit der Zerlegung (fünfte Auslassung). \textbf{Das Lehrbuch der Sensitivitätsanalyse}, aus der Werkstatt, die sie zur Pflicht gemacht hat.}

\lit{27}{Vehtari, A., Gelman, A., Simpson, D., Carpenter, B., Bürkner, P.-C.: \textit{Rank-Normalization, Folding, and Localization: An Improved $\hat R$}. Bayesian Anal.\ 16, 2021.}{Die Weiterentwicklung von $\hat R$ aus dem Forschungskasten in Kapitel~8: robust gegen schwere Ränder, mit getrenntem wirksamem $n$ für Ränder und Mitte. \textbf{Was Stan heute berechnet, steht hier.}}

\lit{28}{Waldmann, K.-H., Helm, W.: \textit{Simulation stochastischer Systeme --- Eine anwendungsorientierte Einführung}. Springer Gabler 2016.}{\textbf{Auf Deutsch}: Zufallszahlen, Erzeugung von Verteilungen, ereignisgesteuerte Simulation von Warteschlangen und Lagersystemen, statistische Auswertung mit Einschwingphase und Blockmitteln --- die Kapitel~1, 2, 4 und 10 dieses Bandes. Mit Beispielen aus Produktion und Logistik, die sich eins zu eins auf die Station übertragen lassen.}

